% TOP_PROBLEM: {"id": 2307034, "title": "Research Problems in Function Theory — Problem 7.34", "category_id": 8, "category": {"id": 8, "name": "analysis", "display_name": "Analysis", "description": "Limits, continuity, calculus, and function theory.", "slug": "analysis", "order_index": 8, "created_at": "2026-07-31T15:26:25.671Z"}, "status": "open"}
% FIRST_POSTED: null
% ENTRY_KIND: statement_only
% Imported from: batch05.tex; UnsolvedMath ID 2307034
\documentclass[11pt]{article}
\usepackage{amsmath,amssymb,mathtools}
\usepackage{fontspec,unicode-math}
\setmainfont{Latin Modern Roman}
\setmathfont{Latin Modern Math}
\usepackage{xurl,hyperref}
\newcommand{\sref}[2]{\hyperlink{source:#1:#2}{[S#2]}}
\newcommand{\cataloguescope}[1]{\paragraph{Mathematical statement.}#1}
\begin{document}
% UnsolvedMath ID 2307034; source catalogue code AMR-022-7034
\setcounter{section}{2307033}
\section{Research Problems in Function Theory — Problem 7.34}\label{2307034}
{\small\sffamily UnsolvedMath ID 2307034 · Analysis}
\subsection{Definitions and mathematical statement}

\paragraph{Shared notation.}$\mathbb N=\{1,2,\ldots\}$, and $\mathbb Z,\mathbb Q,\mathbb R,\mathbb C$ denote the integers, rationals, reals and complex numbers. Any other conventions are specified locally.


Fix $n\in\mathbb N$ and positive real $\beta_1,\ldots,\beta_n$. For arbitrary $\zeta_1,\ldots,\zeta_n\in\mathbb C$, use the branches of $(1-\zeta_jz)^{\beta_j}$ analytic near zero with value one at zero, and expand
\[\prod_{j=1}^n(1-\zeta_jz)^{\beta_j}=1+\sum_{k\ge1}a_kz^k.\]
Restrict to choices for which every $a_k$ is real. Let $N(\beta_1,\ldots,\beta_n)$ be the least positive integer such that every admissible choice of the $\zeta_j$ has some $a_k\le0$ with $1\le k\le N$.

\paragraph{Statement recorded in the supplied source.}
Find the sharp value, or a good upper bound, for the finite integer $N(\beta_1,\ldots,\beta_n)$ in terms of the exponents, uniformly over all admissible $\zeta_j$. \sref{2307034}{2}
The source records finiteness and the bound $N\le\beta_1+\cdots+\beta_n+1$ when all exponents are integers. \sref{2307034}{2}
\subsection{Short English statement}
Bound the first nonpositive Taylor coefficient of the indicated product uniformly in its complex factors, using only the positive exponents.
\subsection{Sources}
\begin{description}
\item[\hypertarget{source:2307034:1}{S1}] ULAM.AI and the UnsolvedMath Contributors, \emph{UnsolvedMath: A Curated Collection of Open Mathematics Problems}, record 2307034 (AMR-022-7034). CC BY4.0. \url{https://huggingface.co/datasets/ulamai/UnsolvedMath}.
\item[\hypertarget{source:2307034:2}{S2}] W. K. Hayman and E. F. Lingham, \emph{Research Problems in Function Theory} (2018), Chapter7, Problem7.34 and its complete update. \url{https://arxiv.org/abs/1809.07200}.
\end{description}
\label{end:2307034}
\end{document}
